% !TEX program = lualatex
% Auto-generated from YAML (v3) — 後期13: データサイエンス・AI（6）：生成AIと倫理的課題

\documentclass[aspectratio=43,professionalfonts,handout]{beamer}
\usepackage{beamer_template}

\newcommand{\LectureNum}{13}
\newcommand{\LectureTitle}{データサイエンス・AI（6）：生成AIと倫理的課題}
\newcommand{\CourseName}{情報リテラシー}
\newcommand{\TextPages}{対応なし}
\newcommand{\NextTopic}{データサイエンス・AI（7）：AIモデルの評価}
\newcommand{\NextPages}{対応なし}

\begin{document}
% --- Slide 1: title ---

\begin{frame}{\CourseName\quad 第\LectureNum 回「\LectureTitle 」}
  {\small \TermName\quad \TeacherName\quad ／\quad 教科書 \TextPages}

  \vfill

  \begin{block}{今日の流れ}
    \begin{enumerate}
      \item 生成AIとは
      \item 生成AIの種類
      \item 生成AIの倫理的課題
      \item NotebookLMとCopilotの比較
      \item 批判的思考の重要性
    \end{enumerate}
  \end{block}
\end{frame}

% --- Slide 2: terms ---

\begin{frame}{生成AIとは}
  \begin{description}[labelwidth=6em]
    \item[\kw{生成AI}]
      新しいコンテンツを生成するAI。テキスト・画像・音声・動画などを生成する。深層学習（第12回）の応用として発展
    \item[\kw{テキスト生成AI}]
      ChatGPT・Claude・Copilot（会話・文章作成・翻訳・要約）
    \item[\kw{画像生成AI}]
      Stable Diffusion・DALL-E（テキストから画像を生成）
    \item[\kw{その他}]
      音声生成（Siri・音声合成）・動画生成・音楽生成・コード生成
  \end{description}
\end{frame}

% --- Slide 3: table ---

\begin{frame}{生成AIの倫理的課題}
  \centering
  \small
  \begin{tabular}{lll}
    \toprule
    課題 & 内容 & 対策 \\
    \midrule
    バイアスと公平性 & 学習データのバイアスが反映される（第9回の採用AIと同じ問題） & 批判的思考・データの多様性確保 \\
    ハルシネーション（誤情報） & AIが事実でない情報をもっともらしく生成する & ファクトチェック・複数情報源の確認 \\
    著作権 & 学習データの著作権・生成物の著作権が誰のものか不明確。法整備が追いついていない & 利用規約の確認・引用明示 \\
    プライバシー & 入力した情報が学習に使われる可能性・個人情報の漏洩リスク & 機密情報を入力しない \\
    依存と思考力低下 & 自分で考えなくなる・批判的思考力・創造性の減少 & AIは補助ツール，最終判断は人間 \\
    悪用 & フェイクニュース・ディープフェイク・なりすまし・試験でのカンニング & 倫理的使用・リテラシー教育 \\
    \bottomrule
  \end{tabular}
\end{frame}

% --- Slide 4: twocol ---

\begin{frame}{NotebookLMとCopilotの比較}
  \begin{columns}[T]
    \begin{column}{0.48\textwidth}
      \begin{block}{NotebookLM（Google提供）}
        \small
        \begin{itemize}
          \item アップロードした資料のみを参照
          \item ドキュメント分析に特化
          \item 引用元が明確
          \item 資料ベースで信頼性が高い
          \item → 復習・資料分析に最適
          \item URL：https://notebooklm.google.com
        \end{itemize}
      \end{block}
    \end{column}
    \begin{column}{0.48\textwidth}
      \begin{exampleblock}{Copilot（Microsoft提供）}
        \small
        \begin{itemize}
          \item GPT-4ベース・汎用的な質問応答
          \item テキスト・画像・コード生成
          \item 一般知識で回答
          \item 引用元が不明確・ハルシネーションの可能性あり
          \item → アイデア出しに最適
          \item URL：https://copilot.microsoft.com
        \end{itemize}
      \end{exampleblock}
    \end{column}
  \end{columns}
\end{frame}

% --- Slide 5: notice ---

\begin{frame}{批判的思考の重要性}
  \begin{noticebox}[批判的思考の重要性]
    \begin{itemize}
      \item AIの回答を鵜呑みにしない：NotebookLMもCopilotも完璧ではない。間違った情報を出すことがある
      \item 批判的に読む：おかしいと思ったら資料を確認・複数の情報源を比較・最終判断は自分で行う
      \item AIは学習を助けるツール：考えることを放棄しない・理解を深めるために使う・自分の言葉でまとめる
      \item 👉 工学技術者として：AIを使いこなしながら，その限界も理解することが重要
    \end{itemize}
  \end{noticebox}
\end{frame}

% --- Slide 6: practice ---

\begin{frame}{実習：NotebookLMとCopilotを使い比べる}
  \begin{block}{手順}
    \begin{enumerate}
      \item Part 1：NotebookLMにアクセス → 第9〜12回ハンドアウト（PDF）をアップロード → 要約を確認 → 質問応答（例：「機械学習の3つの分類は？」）
      \item Part 2：Copilotにアクセス → NotebookLMと同じ質問をする → 回答を比較
      \item Part 3：2つのツールの違いをまとめる
    \end{enumerate}
  \end{block}
\end{frame}

% --- Slide 7: assignment ---

\begin{frame}{課題・次回予告}
  \begin{importantbox}[課題]
    NotebookLMとCopilot比較レポート（800-1000字，Word形式，Teams提出）。構成：①はじめに ②NotebookLMの結果（質問3つ以上） ③Copilotの結果 ④比較と分析 ⑤総合考察。注意：AIの回答をそのまま写さず，自分の言葉でまとめること
  \end{importantbox}

  \vfill

  \begin{block}{次回}
    「\NextTopic 」（教科書 \NextPages ）
  \end{block}

  \vfill

  {\small
  質問があれば授業後またはTeamsで受け付けます。}
\end{frame}

\end{document}
